\documentclass[12pt,a4paper]{report}
\usepackage[footskip=0.2cm]{geometry}
\usepackage[english]{babel}
\usepackage[sfdefault,lf]{carlito}
\usepackage{color}
\usepackage{setspace}
\usepackage{float}
\usepackage{graphicx} %insert figures
\usepackage{fancyhdr}
\usepackage{etoolbox}

\renewcommand{\headrulewidth}{0pt}
\renewcommand{\footrulewidth}{1pt}

\begin{document}
% geometry settings for title page 
    \newgeometry{left=20mm, top=25mm, right=15mm, bottom=35mm}
% The title page settings can be changed in titlepage.tex document
\begin{titlepage}
        \thispagestyle{fancy}
    \begin{figure}[h]
        \vspace{-0.5cm} % changes vertical spacing 
        \centering
        \includegraphics[scale=1.05]{logo_unibo.png}
    \end{figure}

    \vspace{3mm}

    \begin{center}
        \bf {DEPARTMENT OF PHYSICS AND ASTRONOMY ``A. RIGHI''}
    \end{center}

    \begin{center}
        \bf \large{SECOND CYCLE DEGREE}

        \bf \large{PHYSICS}
    \end{center}

    \vspace{15mm}

    \begin{center}
        %
        % INSERT THE TITLE OF YOUR THESIS
        %
        {\huge{\bf 3D U-NET DOMAIN GENERALIZATION IN}}
        \vspace{15pt}\\
        {\huge{\bf FETAL BRAIN MRI SEGMENTATION}}
    \end{center}

    \vspace{50mm} \par \noindent

    \begin{minipage}[t]{0.47\textwidth}
        %
        % INSERT THE NAME OF THE SUPERVISOR WITH ITS TITLE (DR. OR PROF.)
        %
        \vspace*{-19pt}
        {\large{\bf Supervisor
        \vspace{2mm}\\
        Prof.\ Daniel Remondini\\\\
        %
        % INSERT THE NAME OF THE CO-SUPERVISOR WITH ITS TITLE (DR. OR PROF.)
        %
        % IF THERE ARE NO CO-SUPERVISORS REMOVE THE FOLLOWING 5 LINES
        %
        Co-supervisors
        \vspace{2mm}\\
        Dr.\ Nico Curti\\
        Dr.\ Gerard Martí Juan\\\\}}
    \end{minipage}
    %
    \hfill
    %
    \begin{minipage}[L]{0.23\textwidth} \textcolor{black}{
        {\large{\bf  Defended by
        \vspace{2mm}\\
        \noindent{Simone Chiarella}}}
        }
    \end{minipage}
    %
    \cfoot[C]{
        \\
        \bf 18--19/12/2025\\
        \vspace{3mm}
        \bf Academic Year 2024/2025
    }

\end{titlepage}
\restoregeometry

\newpage\null\thispagestyle{empty}\newpage

\renewcommand\abstractname{\textsc{Abstract}}
{
  \fontfamily{bch}\selectfont\begin{abstract}
    This work investigates domain generalization of a 3D U-Net (nnU-Net v2.4.1) for fetal brain MRI tissue segmentation. The central question was whether GIN-IPA, a promising but not widely studied image augmentation method, would improve the network robustness across acquisition shifts. Three datasets with different super-resolution reconstruction (SRR) were employed: Kispi-mial (small size, both healthy and pathological cases), Kispi-irtk (similar to Kispi-mial but higher quality), and dHCP (larger size, high quality, healthy only). Label sets were harmonized to a common 7-tissue scheme.

    Three strategies were compared within the nnU-Net training loop: (i) default nnU-Net augmentation, comprising standard geometric and photometric transforms; (ii) GIN-IPA, an appearance perturbation that aims to mimic acquisition shifts; and (iii) the combination of both. Models were trained separately on each dataset and assessed on unseen samples from both the source dataset and external datasets, using typical metrics such as the Dice score and others.

    Two conclusions were drawn. First, generalization was primarily determined by data quality and scale: training on dHCP yielded the most stable cross-domain performance, almost independent of augmentation. Second, GIN-IPA was beneficial when source data were limited and differed from the target: training on Kispi-irtk and inferring on dHCP resulted in significant improvements compared to default nnU-Net augmentation. By contrast, stacking the two augmentations was not additive and, at times, detrimental.

    Limitations include the small size of public clinical datasets, the need for label-set harmonization, and hardware constraints that prevented a full exploration of the GIN-IPA potential. The results suggest that multi-center, high-quality fetal MRI with standardized SRR should be prioritized. Within constrained single-source regimes, GIN-IPA may represent a pragmatic option for domain generalization, despite the need for further validation.
  \end{abstract}
}

\end{document}

%This work investigates domain generalization of a 3D U-Net (nnU-Net v2.4.1) for fetal brain MRI tissue segmentation. The central question was whether GIN-IPA, a promising though not broadly studied image augmentation method, would improve the network robustness across acquisition sites and super-resolution reconstruction (SRR) pipelines. Three datasets with distinct characteristics were employed: Kispi-mial (\textsc{mialsrtk} SRR; small size, both neurotypical and pathological cases; $1.5$/$3$\,T,), Kispi-irtk (\textsc{irtk} SRR; same features as Kispi-mial but higher quality), and dHCP (dedicated reconstruction; larger size, high quality, healthy-only atlas; 3\,T). Label sets were harmonized to a common 7-tissue scheme. ``In-domain'' (ID) denoted evaluation on the training dataset; ``out-of-domain'' (OOD) denoted evaluation on unseen datasets with different characteristics.

%Three strategies were compared within the nnU-Net training loop: (i) default nnU-Net augmentation, comprising standard geometric and photometric transforms; (ii) GIN-IPA, an appearance perturbation realized via random convolutions to mimic acquisition shifts; and (iii) the combination of both. Models were trained separately on each dataset and assessed ID and OOD using Dice score, volume similarity, and Hausdorff distance.

%Two conclusions were drawn. First, generalization was dominated by data quality and scale: when trained on dHCP, the most stable cross-domain performance was obtained, almost independent of the augmentation. Second, GIN-IPA was beneficial when training data were limited and differed from the target: with training on Kispi-irtk and inference on dHCP, significant improvements were observed, with respect to default nnU-Net augmentation. Instead, stacking the two augmentations was not additive and, at times, detrimental.

%Limitations include the small size of public clinical datasets, the need for label-set harmonization, and hardware constraints that did not allow for a complete exploration of the GIN-IPA potential. The results suggest that multi-center, high-quality fetal MRI with standardized SRR should be prioritized. Within constrained single-source regimes, GIN-IPA may constitute a pragmatic option for domain generalization, despite the need for further validation.